\section{Introduction}
\label{sec:intro}

The proliferation of deep neural networks in modern applications has led to a paradigm shift in deploying specialized capabilities on edge hardware. Typically, a general-purpose base network (e.g., a pre-trained Vision Transformer \cite{vaswani2017attention} or CLIP \cite{radford2021learning}) is fine-tuned on multiple downstream datasets to construct task-specific experts. Under standard edge paradigms, hosting an independent expert for every task is highly impractical, violating strict physical constraints on flash memory, storage, and thermal limits on low-power smart sensors \cite{gholami2021survey, jacob2018quantization}.

To circumvent this bottleneck, the community has embraced \emph{multi-task model merging} \cite{ilharco2022editing, wortsman2022model, jin2023regmean}, combining multiple single-task experts (sharing a base initialization) into a unified network directly in parameter space. This bypasses the substantial computation and data requirements of joint multi-task training. Among these, Task Arithmetic \cite{ilharco2022editing} offers an elegant solution: by adding task-specific weight differences (task vectors) to the base model, a single consolidated network performs multiple tasks.

However, for physical deployment, model merging in its raw, full-precision form (FP32/FP16) remains incomplete. Under strict hardware constraints, post-training quantization (PTQ) to low-bit formats, such as 8-bit (INT8) or 4-bit (INT4) integers, is an absolute operational necessity. PTQ reduces on-device storage footprints by 50\% to 75\% and alleviates memory bandwidth bottlenecks—the primary drivers of energy consumption and latency on the edge.

Yet, when practitioners attempt to unify model merging and low-bit quantization, they run into a challenging dilemma (as shown in Figure \ref{fig:accuracy_comparison}):
\begin{enumerate}
    \item \textbf{Merge-then-Quantize (M-then-Q):} Merging experts in full precision and subsequently quantizing the unified weights. This post-hoc quantization introduces severe quantization noise that breaks representation alignment and significantly degrades multi-task test accuracy.
    \item \textbf{Quantize-then-Merge (Q-then-M):} Merging experts that have already been quantized to integer formats. This is mathematically invalid; independent integer quantization breaks the continuous linear mode connectivity (LMC) \cite{ainsworth2022git, entezari2021role} between the experts, resulting in catastrophic representational collapse.
\end{enumerate}

To bridge this gap, emerging techniques like \emph{Quantization-Aware Merging} (\texttt{Q-Merge}) and test-time adaptation (TTA) frameworks (such as \texttt{AdaMerging} \cite{yang2024adamerging}) attempt to optimize layer-wise merging coefficients directly on the device. At test time, a stream of unlabeled calibration images from the target domains is routed through the network. The merging coefficients are adjusted to minimize the entropy of the output logit distributions, thereby aligning the parameters under the physical quantization operator.

\begin{figure}[t]
\centering
\includegraphics[width=0.48\textwidth]{results/accuracy_comparison.png}
\caption{\textbf{Average Multi-Task Accuracy Comparison} under 8-Bit post-training quantization on our Vision Transformer (ViT-Tiny) multi-task benchmark. While naive Merge-then-Quantize (M-then-Q) and unconstrained Q-Merge suffer from representation misalignment and overfitting, our proposed \textbf{Q-PolyMerge} (Adam STE) recovers the performance ceiling of unquantized models.}
\label{fig:accuracy_comparison}
\end{figure}

However, adaptive frameworks suffer from a severe vulnerability we term the \textbf{Overfitting-Optimizer Paradox}. On edge devices, test-time calibration streams are necessarily small and transient (e.g., 16 unlabeled images). When an unconstrained optimizer (e.g., Adam \cite{bengio2013estimating} or evolutionary search \cite{akiba2024evolutionary}) searches for independent layer-wise coefficients across $L$ layers and $K$ tasks (56 parameters in our ViT-Tiny setup), it easily overfits to transductive noise, yielding jagged, physically nonsensical trajectories that fail to generalize.

To resolve this, we propose \textbf{Q-PolyMerge}, a hybrid-precision (weight-quantized, activation-float) framework for low-bit model merging that is highly parameter-efficient and robust. Q-PolyMerge projects continuous layer coefficients onto a low-degree polynomial subspace of normalized layer depth. Rather than optimizing independent layer-wise parameters, the optimizer adjusts a small set of polynomial coefficients of degree $d$ (typically $d=2$), achieving several crucial advantages:
\begin{itemize}
    \item \textbf{Extreme Parameter Efficiency:} We reduce the parameter search space by \textbf{78.6\%} (from 56 to 12 parameters).
    \item \textbf{Mathematical Low-Pass Filtering:} Bounding the trajectories to a smooth polynomial acts as a low-pass filter, mathematically removing optimization noise and preventing degenerate entropy-collapse.
    \item \textbf{On-Device Zero-Order Viability:} Our low dimensionality enables highly efficient derivative-free optimization (1+1 ES). Treating the network as a black-box oracle bypasses backpropagation and activation caching, eliminating gradient storage overhead.
    \item \textbf{Integer-Weight Pipeline (with FP Activations):} Storing 100\% of parameters as low-bit integers removes floating-point variables from on-device storage. Crucially, activations and normalization calculations are kept in floating-point during inference to ensure capacity and numerical stability.
\end{itemize}

Through systematic evaluation of a Vision Transformer on MNIST, FashionMNIST, CIFAR-10, and SVHN across three random seeds under 8-bit and 4-bit quantization, we demonstrate that Q-PolyMerge consistently achieves a superior balance of high performance and stability. Under 4-bit PTQ, Q-PolyMerge (Adam STE) achieves \textbf{48.87\%} average accuracy, strictly outperforming unconstrained Q-Merge (\textbf{46.02\%}) and naive post-merge quantization (\textbf{42.92\%}). Under 8-bit PTQ, Q-PolyMerge achieves \textbf{59.76\%} average accuracy, which outperforms naive post-merge (\textbf{55.11\%}) and matches unconstrained Q-Merge (\textbf{60.03\%}) while reducing standard deviation by over 47\% and providing smooth, physically stable layer trajectories.

\textbf{Distinguishing On-Device TTA from Offline Ceilings:} We address an important baseline anomaly: offline, server-side optimization in full precision followed by post-hoc quantization (e.g., standard \texttt{AdaMerging (Adam $\to$ Low-Bit)}) achieves high accuracies (e.g., \textbf{62.27\%} in 8-bit and \textbf{50.20\%} in 4-bit), slightly outperforming our direct quantization-aware on-device methods. However, we clarify that such offline, first-order methods are physically impossible to run on physical low-power edge nodes for test-time adaptation. Performing first-order backpropagation requires storing weights in full precision and caching intermediate activation tensors, consuming an enormous SRAM footprint (approx. \textbf{158.4 MB} for ViT-Tiny) that completely exceeds the typical microcontroller internal SRAM budget ($\le 2$ MB). Under the physically viable zero-order regime (which only runs forward passes and requires no activation caching), unconstrained search collapses due to the curse of dimensionality, while Q-PolyMerge successfully filters out optimization noise, delivering substantial absolute improvements on physical edge devices.